%%%%%%%%%%%%%%%%
%%% exod 26 %%%
%%%%%%%%%%%%%%%%

\Chapter{26}

\Verse{1}
{
	\hP{וְאֶת הַמִּשְׁכָּן תַּעֲשֶׂה עֶשֶׂר יְרִיעֹת שֵׁשׁ משְׁזָר וּתְכֵלֶת וְאַרְגָּמָן וְתֹלַעַת שָׁנִי כְּרֻבִים מַעֲשֵׂה חֹשֵׁב תַּעֲשֶׂה אֹתָם׃}
}
{
	\eP{
		“And you shall make the Tabernacle with ten curtains of
		woven linen and blue and purple and scarlet, with
		cherubs—you shall make them designer’s work.
	}
}
{
%	\footnote{
%		make the Tabernacle. No one has ever figured out how the Tabernacle is
%		put together. The components are given here, but the text does not say
%		how they are put together. Twenty frames are used for each side wall and
%		six for the rear wall (plus two more for supporting the rear corners).
%		The frames are ten cubits tall and “a cubit and a half cubit” wide.
%		People have therefore assumed that the Tabernacle is thirty cubits long
%		by ten cubits wide, and ten cubits high. Curtains fit over these frames,
%		and the curtains are made of two giant fabrics. Each fabric is made of
%		five pieces sewn together; the pieces are four cubits by twenty-eight
%		cubits wide. So the two fabrics are each twenty by twenty-eight. They
%		are attached together with gold clasps, making a double fabric that is
%		forty by twenty-eight. When this forty-cubit piece is placed over the
%		thirty-cubit-long structure of frames, it covers the top and has ten
%		cubits left over to cover the ten-cubit-high rear wall. In support of
%		this, people have noted that Solomon’s Temple is sixty cubits long by
%		twenty wide, so the Temple and Tabernacle are in two-to-one proportion.
%		But there are many things wrong with this. (1) The Temple is thirty
%		cubits high, so the proportion there is three-to-one, not two-to-one.
%		(2) The ten-cubit width is wrong, because six frames of a cubit and a
%		half each make only a nine-cubit width. (3) The twenty-eight-cubits-wide
%		curtains do not fit this structure. Its two sidewalls are each ten
%		cubits high, and its ceiling is ten cubits across, which requires a
%		cover that is thirty cubits wide. The twenty-eight-cubits-wide curtain
%		leaves a whole cubit (eighteen inches) uncovered on each side. (4) A
%		second set of curtains is made into two large fabrics like the first
%		set, but this second set (made of goats’ hair) is made of curtains that
%		are four cubits by thirty, instead of four by twenty-eight, and there is
%		an extra, eleventh curtain at one end. The extra curtain has no purpose,
%		and the thirty-cubits-length hangs down past the bottom of the inner
%		fabric. (5) The gold clasps cannot be seen in this arrangement. Also,
%		this arrangement leaves questions: (6) Why would the frames be a cubit
%		and a half wide instead of just one cubit or two cubits? And why does it
%		say “a cubit and a half cubit” instead of the normal wording, “a cubit
%		and a half”? (7) Why are the curtains in groups of five, which are each
%		four cubits wide? What is significant about four cubits? All of these
%		problems and questions point to a different way that the Tabernacle is
%		put together. I propose that the reason that the frames are “a cubit and
%		a half cubit” is because they overlap each other by the extra
%		half-cubit. This is more stable and better for ventilation than if they
%		stood side by side (Figures 1 and 2). The twenty-frame length by
%		six-frame width means a structure that is twenty cubits by six cubits on
%		the inside. If the frames are a half-cubit wide, then it is twenty by
%		eight on the outside (Figure 3). Now, if we take the two large fabrics
%		that are twenty by twenty-eight cubits, and we connect them with the
%		clasps, and then we fold it in half along the clasps, it fits this
%		structure. The gold clasps now are visible: they surround the entrance
%		to the Tabernacle. The twenty-cubit width of the fabric matches the
%		twenty-cubit length of the Tabernacle. And the twenty-eight-cubit length
%		of the fabric goes up one side of the Tabernacle (ten cubits high),
%		across the ceiling (eight cubits), and down the other side (ten cubits).
%		A temple has been excavated in Israel at Arad. Its measurements are
%		twenty cubits by six cubits, the same as the (inside dimensions of the)
%		Tabernacle. On the significance of these dimensions, see the comment on
%		26:30.
%	}
}

\Verse{2}
{
	\hP{אֹרֶךְ הַיְרִיעָה הָאַחַת שְׁמֹנֶה וְעֶשְׂרִים בָּאַמָּה וְרֹחַב אַרְבַּע בָּאַמָּה הַיְרִיעָה הָאֶחָת מִדָּה אַחַת לְכל הַיְרִיעֹת׃}
}
{
	\eP{
		The length of one curtain shall be twenty-eight in cubits,
		and the width of one curtain four in cubits: one size to
		all the curtains.
	}
}
{
}

\Verse{3}
{
	\hP{חֲמֵשׁ הַיְרִיעֹת תִּהְיֶיןָ חֹבְרֹת אִשָּׁה אֶל אֲחֹתָהּ וְחָמֵשׁ יְרִיעֹת חֹבְרֹת אִשָּׁה אֶל אֲחֹתָהּ׃}
}
{
	\eP{
		Five curtains will be connected, each to its sister-piece,
		and five curtains connected, each to its sister-piece.
	}
}
{
}

\Verse{4}
{
	\hP{וְעָשִׂיתָ לֻלְאֹת תְּכֵלֶת עַל שְׂפַת הַיְרִיעָה הָאֶחָת מִקָּצָה בַּחֹבָרֶת וְכֵן תַּעֲשֶׂה בִּשְׂפַת הַיְרִיעָה הַקִּיצוֹנָה בַּמַּחְבֶּרֶת הַשֵּׁנִית׃}
}
{
	\eP{
		And you shall make loops of blue on the side of the one
		curtain at the end of the connected group, and you shall
		do so in the side of the end curtain in the second
		connected group.
	}
}
{
}

\Verse{5}
{
	\hP{חֲמִשִּׁים לֻלָאֹת תַּעֲשֶׂה בַּיְרִיעָה הָאֶחָת וַחֲמִשִּׁים לֻלָאֹת תַּעֲשֶׂה בִּקְצֵה הַיְרִיעָה אֲשֶׁר בַּמַּחְבֶּרֶת הַשֵּׁנִית מַקְבִּילֹת הַלֻּלָאֹת אִשָּׁה אֶל אֲחֹתָהּ׃}
}
{
	\eP{
		You shall make fifty loops in the end of the curtain that
		is in the second connected group, with the loops parallel:
		each to its sister-piece.
	}
}
{
}

\Verse{6}
{
	\hP{וְעָשִׂיתָ חֲמִשִּׁים קַרְסֵי זָהָב וְחִבַּרְתָּ אֶת הַיְרִיעֹת אִשָּׁה אֶל אֲחֹתָהּ בַּקְּרָסִים וְהָיָה הַמִּשְׁכָּן אֶחָד׃}
}
{
	\eP{
		And you shall make fifty clasps of gold and connect the
		curtains, each to its sister-piece, with the clasps. And
		the Tabernacle will be one.
	}
}
{
%	\footnote{
%		the Tabernacle will be one. In the simplest sense, this means that once
%		they sew all the curtains together and connect all the loops and clasps,
%		the Tabernacle will be one whole piece. But the wording—“the Tabernacle
%		will be one”—also fits with the centrality of the Tabernacle to Israel’s
%		monotheism. One God, one Tabernacle, one altar, only one place of
%		worship. Later it will be confirmed by divine commandment that Israel
%		may have only one place of sacrifice, and that this one place is in
%		front of the Tabernacle (Leviticus 17). There cannot be more than one
%		place to worship God, because that might suggest that there is more than
%		one God. Or, to put it more essentially: it is inconceivable in the
%		first place in this new religion of a single deity that there could
%		possibly be more than a single sanctuary. Those who claim that Israel’s
%		religion was not a real monotheism until a very late stage (Second
%		Temple) have not appreciated the significance of the commandment of
%		centralization of worship at a single location.
%	}
}

\Verse{7}
{
	\hP{וְעָשִׂיתָ יְרִיעֹת עִזִּים לְאֹהֶל עַל הַמִּשְׁכָּן עַשְׁתֵּי עֶשְׂרֵה יְרִיעֹת תַּעֲשֶׂה אֹתָם׃}
}
{
	\eP{
		“And you shall make curtains of goats’ hair for a tent
		over the Tabernacle. You shall make them eleven curtains.
	}
}
{
}

\Verse{8}
{
	\hP{אֹרֶךְ הַיְרִיעָה הָאַחַת שְׁלֹשִׁים בָּאַמָּה וְרֹחַב אַרְבַּע בָּאַמָּה הַיְרִיעָה הָאֶחָת מִדָּה אַחַת לְעַשְׁתֵּי עֶשְׂרֵה יְרִיעֹת׃}
}
{
	\eP{
		The length of one curtain shall be thirty in cubits, and
		the width of one curtain four in cubits: one size for
		eleven curtains.
	}
}
{
}

\Verse{9}
{
	\hP{וְחִבַּרְתָּ אֶת חֲמֵשׁ הַיְרִיעֹת לְבָד וְאֶת שֵׁשׁ הַיְרִיעֹת לְבָד וְכָפַלְתָּ אֶת הַיְרִיעָה הַשִּׁשִּׁית אֶל מוּל פְּנֵי הָאֹהֶל׃}
}
{
	\eP{
		And you shall connect five of the curtains by themselves
		and six of the curtains by themselves. And you shall
		double-fold the sixth curtain opposite the front of the
		tent.
	}
}
{
%	\footnote{
%		double-fold the sixth curtain. The extra, eleventh curtain of the tent
%		is four cubits wide. When folded, each half of it covers half of the
%		eight-cubits-wide back of the Tabernacle (Figure 4). This also answers
%		the question of why the curtains are four cubits wide. And it further
%		confirms that the Tabernacle is eight cubits wide. Footnote 26:9.
%		opposite the front of the tent. Meaning: it falls along the back wall,
%		facing the front. Thus v. 12 refers to “the hanging one that is left
%		over among the tent’s curtains” and instructs that “you shall hang half
%		of the leftover curtain on the back parts of the Tabernacle.”
%	}
}

\Verse{10}
{
	\hP{וְעָשִׂיתָ חֲמִשִּׁים לֻלָאֹת עַל שְׂפַת הַיְרִיעָה הָאֶחָת הַקִּיצֹנָה בַּחֹבָרֶת וַחֲמִשִּׁים לֻלָאֹת עַל שְׂפַת הַיְרִיעָה הַחֹבֶרֶת הַשֵּׁנִית׃}
}
{
	\eP{
		And you shall make fifty loops on the side of the one end
		curtain in the connected group and fifty loops on the side
		of the curtain of the second connected group.
	}
}
{
}

\Verse{11}
{
	\hP{וְעָשִׂיתָ קַרְסֵי נְחֹשֶׁת חֲמִשִּׁים וְהֵבֵאתָ אֶת הַקְּרָסִים בַּלֻּלָאֹת וְחִבַּרְתָּ אֶת הָאֹהֶל וְהָיָה אֶחָד׃}
}
{
	\eP{
		And you shall make fifty clasps of bronze and bring the
		clasps into the loops and connect the tent. And it will be
		one.
	}
}
{
}

\Verse{12}
{
	\hP{וְסֶרַח הָעֹדֵף בִּירִיעֹת הָאֹהֶל חֲצִי הַיְרִיעָה הָעֹדֶפֶת תִּסְרַח עַל אֲחֹרֵי הַמִּשְׁכָּן׃}
}
{
	\eP{
		And the hanging one that is left over among the tent’s
		curtains: you shall hang half of the leftover curtain on
		the back parts of the Tabernacle.
	}
}
{
}

\Verse{13}
{
	\hP{וְהָאַמָּה מִזֶּה וְהָאַמָּה מִזֶּה בָּעֹדֵף בְּאֹרֶךְ יְרִיעֹת הָאֹהֶל יִהְיֶה סָרוּחַ עַל צִדֵּי הַמִּשְׁכָּן מִזֶּה וּמִזֶּה לְכַסֹּתוֹ׃}
}
{
	\eP{
		And the cubit on this side and the cubit on that side in
		the leftover in the length of the tent’s curtains shall be
		hung on the sides of the Tabernacle, on this side and on
		that side, to cover it.
	}
}
{
%	\footnote{
%		the cubit on this side and the cubit on that side in the leftover in the
%		length. This second set of curtains is thirty cubits long. When it is
%		spread over the inside set of curtains, which is only twenty-eight
%		cubits, there is thus a cubit left over on each side. It “shall be hung
%		on the sides of the Tabernacle, on this side and on that side, to cover
%		it.” That is, it is folded under the inside linen fabric so as to
%		protect it from touching the ground.
%	}
}

\Verse{14}
{
	\hP{וְעָשִׂיתָ מִכְסֶה לָאֹהֶל עֹרֹת אֵילִם מְאדָּמִים וּמִכְסֵה עֹרֹת תְּחָשִׁים מִלְמָעְלָה׃}
}
{
	\eP{
		And you shall make a covering for the tent of rams’ skins
		dyed red and a covering of leather skins above.
	}
}
{
}

\Verse{15}
{
	\hP{וְעָשִׂיתָ אֶת הַקְּרָשִׁים לַמִּשְׁכָּן עֲצֵי שִׁטִּים עֹמְדִים׃}
}
{
	\eP{
		“And you shall make the frames for the Tabernacle of
		acacia wood, standing.
	}
}
{
%	\footnote{
%		frames. Trellises—not solid boards or planks. If they were solid, then
%		the linen curtains would be sandwiched between the boards on the inside
%		and the goats’ hair and leather coverings on the outside and could never
%		be seen.
%	}
}

\Verse{16}
{
	\hP{עֶשֶׂר אַמּוֹת אֹרֶךְ הַקָּרֶשׁ וְאַמָּה וַחֲצִי הָאַמָּה רֹחַב הַקֶּרֶשׁ הָאֶחָד׃}
}
{
	\eP{
		The frame’s length shall be ten cubits, and the width of
		one frame a cubit and a half cubit.
	}
}
{
}

\Verse{17}
{
	\hP{שְׁתֵּי יָדוֹת לַקֶּרֶשׁ הָאֶחָד מְשֻׁלָּבֹת אִשָּׁה אֶל אֲחֹתָהּ כֵּן תַּעֲשֶׂה לְכֹל קַרְשֵׁי הַמִּשְׁכָּן׃}
}
{
	\eP{
		One frame has two projections, each aligned with its
		sister-piece: so you shall make for all the Tabernacle’s
		frames.
	}
}
{
}

\Verse{18}
{
	\hP{וְעָשִׂיתָ אֶת הַקְּרָשִׁים לַמִּשְׁכָּן עֶשְׂרִים קֶרֶשׁ לִפְאַת נֶגְבָּה תֵימָנָה׃}
}
{
	\eP{
		And you shall make the frames for the Tabernacle: twenty
		frames for the south side—to the south.
	}
}
{
}

\Verse{19}
{
	\hP{וְאַרְבָּעִים אַדְנֵי כֶסֶף תַּעֲשֶׂה תַּחַת עֶשְׂרִים הַקָּרֶשׁ שְׁנֵי אֲדָנִים תַּחַת הַקֶּרֶשׁ הָאֶחָד לִשְׁתֵּי יְדֹתָיו וּשְׁנֵי אֲדָנִים תַּחַת הַקֶּרֶשׁ הָאֶחָד לִשְׁתֵּי יְדֹתָיו׃}
}
{
	\eP{
		And you shall make forty bases of silver under the twenty
		frames: two bases under one frame for its two projections
		and two bases under each other frame for its two
		projections.
	}
}
{
}

\Verse{20}
{
	\hP{וּלְצֶלַע הַמִּשְׁכָּן הַשֵּׁנִית לִפְאַת צָפוֹן עֶשְׂרִים קָרֶשׁ׃}
}
{
	\eP{
		And for the Tabernacle’s second side, for the north side:
		twenty frames
	}
}
{
}

\Verse{21}
{
	\hP{וְאַרְבָּעִים אַדְנֵיהֶם כָּסֶף שְׁנֵי אֲדָנִים תַּחַת הַקֶּרֶשׁ הָאֶחָד וּשְׁנֵי אֲדָנִים תַּחַת הַקֶּרֶשׁ הָאֶחָד׃}
}
{
	\eP{
		and their forty bases of silver, two bases under one frame
		and two bases under each other frame.
	}
}
{
}

\Verse{22}
{
	\hP{וּלְיַרְכְּתֵי הַמִּשְׁכָּן יָמָּה תַּעֲשֶׂה שִׁשָּׁה קְרָשִׁים׃}
}
{
	\eP{
		And for the rear of the Tabernacle, to the west, you shall
		make six frames.
	}
}
{
}

\Verse{23}
{
	\hP{וּשְׁנֵי קְרָשִׁים תַּעֲשֶׂה לִמְקֻצְעֹת הַמִּשְׁכָּן בַּיַּרְכָתָיִם׃}
}
{
	\eP{
		And you shall make two frames for the Tabernacle’s corners
		in the rear,
	}
}
{
}

\Verse{24}
{
	\hP{וְיִהְיוּ תֹאֲמִם מִלְּמַטָּה וְיַחְדָּו יִהְיוּ תַמִּים עַל רֹאשׁוֹ אֶל הַטַּבַּעַת הָאֶחָת כֵּן יִהְיֶה לִשְׁנֵיהֶם לִשְׁנֵי הַמִּקְצֹעֹת יִהְיוּ׃}
}
{
	\eP{
		and they shall be doubles from below, and together they
		shall be integrated on its top to one ring. It shall be so
		for the two of them; they will be for the two corners.
	}
}
{
}

\Verse{25}
{
	\hP{וְהָיוּ שְׁמֹנָה קְרָשִׁים וְאַדְנֵיהֶם כֶּסֶף שִׁשָּׁה עָשָׂר אֲדָנִים שְׁנֵי אֲדָנִים תַּחַת הַקֶּרֶשׁ הָאֶחָד וּשְׁנֵי אֲדָנִים תַּחַת הַקֶּרֶשׁ הָאֶחָד׃}
}
{
	\eP{
		And they shall be eight frames and their bases of silver,
		sixteen bases, two bases under one frame and two bases
		under each other frame.
	}
}
{
}

\Verse{26}
{
	\hP{וְעָשִׂיתָ בְרִיחִם עֲצֵי שִׁטִּים חֲמִשָּׁה לְקַרְשֵׁי צֶלַע הַמִּשְׁכָּן הָאֶחָד׃}
}
{
	\eP{
		And you shall make bars of acacia wood: five for the
		frames of one side of the Tabernacle
	}
}
{
}

\Verse{27}
{
	\hP{וַחֲמִשָּׁה בְרִיחִם לְקַרְשֵׁי צֶלַע הַמִּשְׁכָּן הַשֵּׁנִית וַחֲמִשָּׁה בְרִיחִם לְקַרְשֵׁי צֶלַע הַמִּשְׁכָּן לַיַּרְכָתַיִם יָמָּה׃}
}
{
	\eP{
		and five bars for the frames of the second side of the
		Tabernacle and five bars for the frames of the side of the
		Tabernacle at the rear, to the west,
	}
}
{
}

\Verse{28}
{
	\hP{וְהַבְּרִיחַ הַתִּיכֹן בְּתוֹךְ הַקְּרָשִׁים מַבְרִחַ מִן הַקָּצֶה אֶל הַקָּצֶה׃}
}
{
	\eP{
		and the middle bar through the frames extending from end
		to end.
	}
}
{
}

\Verse{29}
{
	\hP{וְאֶת הַקְּרָשִׁים תְּצַפֶּה זָהָב וְאֶת טַבְּעֹתֵיהֶם תַּעֲשֶׂה זָהָב בָּתִּים לַבְּרִיחִם וְצִפִּיתָ אֶת הַבְּרִיחִם זָהָב׃}
}
{
	\eP{
		And you shall plate the frames with gold, and you shall
		make their rings gold, housings for the bars, and you
		shall plate the bars with gold.
	}
}
{
}

\Verse{30}
{
	\hP{וַהֲקֵמֹתָ אֶת הַמִּשְׁכָּן כְּמִשְׁפָּטוֹ אֲשֶׁר הרְאֵיתָ בָּהָר׃}
}
{
	\eP{
		And you shall set up the Tabernacle according to the model
		of it that you were shown in the mountain.
	}
}
{
%	\footnote{
%		set up the Tabernacle. What difference do these cubits make? Why does it
%		matter whether the Tabernacle is eight by twenty cubits or ten by
%		thirty? Because eight by twenty cubits, and ten cubits high, is the size
%		of the space under the wings of the cherubs inside the Holy of Holies in
%		the Temple in Jerusalem. Either the Tabernacle really stood in that
%		space, or the space represented it symbolically while the Tabernacle
%		itself was stored away somewhere inside the building. The Tanak reports
%		that the Tabernacle is in fact brought to the Temple on the day of its
%		dedication (1 Kings 8:4; 2 Chr 5:5). The Talmud reports that it was in
%		fact thus stored beneath the Temple (Soa 9a). The ancient historian
%		Josephus says that the effect of the cherubs’ wings in the Temple was to
%		appear like a tent (Ant. 8.103). A psalm says: “I’ll reside in your tent
%		forever, I’ll conceal myself in the hidden place of your wings” (Ps
%		61:5). The book of 1 Chronicles (9:23) refers to the Temple as “the
%		House of the Tent” (bêt h ’hel). It also speaks of the “Tabernacle of
%		the House of God” (6:33; see also 1 Chr 23:22; 2 Chr 29:6–7; 24:6). What
%		this means is that the Temple was not just the successor to the
%		Tabernacle. The Tabernacle, revered as the ancient channel to God, was
%		actually located inside the Temple. So all the laws in the Torah that
%		command that something be done at the Tabernacle applied as long as the
%		first Temple stood in Jerusalem. By the time of the Temple’s destruction
%		(587 B.C.E.), the Tabernacle and its contents—the ark, the
%		tablets—disappear from the Bible. The rabbinic tradition is that the
%		presence of God (Shechinah; which, by the way, has the same root as the
%		word for the Tabernacle: miškn) was in the first Temple but not in the
%		second (Yoma 21b; Rashi on Gen 9:27). This contributes to an impression
%		of growing distance from the original feeling of a direct connection
%		with God in the Hebrew Bible. This construction of the Tabernacle also
%		goes against the dominant view of biblical scholars for the past
%		century: that the Tabernacle was not real, that it was simply a fiction,
%		invented to stand for the second Temple in Jerusalem. That view was
%		based on the two-to-one correspondence of the measurements of the
%		Tabernacle and the Temple. In the first place, we have seen that the
%		correspondence is not correct. (See point number one in the comment on
%		26:1 above.) Second, those measurements were based on the first Temple,
%		not the second, in any case. And, third, this chapter is part of a text
%		that is written in Classical Hebrew, which comes from a period long
%		before the second Temple.
%	}
}

\Verse{31}
{
	\hP{וְעָשִׂיתָ פָרֹכֶת תְּכֵלֶת וְאַרְגָּמָן וְתוֹלַעַת שָׁנִי וְשֵׁשׁ משְׁזָר מַעֲשֵׂה חֹשֵׁב יַעֲשֶׂה אֹתָהּ כְּרֻבִים׃}
}
{
	\eP{
		“And you shall make a pavilion of blue and purple and
		scarlet and woven linen; he shall make them designer’s
		work with cherubs.
	}
}
{
%	\footnote{
%		pavilion. Hebrew prket. In the second Temple this came to be a curtain,
%		veiling the Holy of Holies. Subsequently in synagogues it likewise has
%		been a curtain, veiling the ark. But in the Tabernacle it is a pavilion,
%		a tent within the tent, not a veil. It hangs on four posts, it provides
%		a cover over the ark (wsakkt ‘al h’rn; Exod 40:3), it is over the
%		Testimony (Exod 27:21), and it is referred to as “the covering
%		pavilion,” prket hammsk (Exod 40:21; Num 4:5). (It is the same root word
%		that is used earlier to describe how the cherubs’ wings “cover over” the
%		atonement dais; 25:20.) Such a sukkh is referred to in Ps 27:5 and in
%		Lam 2:6 as well. Made of fabric identical to that of the Tabernacle, it
%		provides the Tabernacle’s necessary rear fabric wall. In the Talmudic
%		age, when the prket had come to be a curtain, the sages assumed that it
%		had always been a curtain, and so they were puzzled that in the Torah it
%		is pictured as a sukkh and that it is described as being “over” the ark
%		(Sukkah 7b; cf. Soa 37a; Menahot 62a, 98a).
%	}
}

\Verse{32}
{
	\hP{וְנָתַתָּה אֹתָהּ עַל אַרְבָּעָה עַמּוּדֵי שִׁטִּים מְצֻפִּים זָהָב וָוֵיהֶם זָהָב עַל אַרְבָּעָה אַדְנֵי כָסֶף׃}
}
{
	\eP{
		And you shall place it on four acacia columns plated with
		gold, their hooks of gold, on four bases of silver.
	}
}
{
}

\Verse{33}
{
	\hP{וְנָתַתָּה אֶת הַפָּרֹכֶת תַּחַת הַקְּרָסִים וְהֵבֵאתָ שָׁמָּה מִבֵּית לַפָּרֹכֶת אֵת אֲרוֹן הָעֵדוּת וְהִבְדִּילָה הַפָּרֹכֶת לָכֶם בֵּין הַקֹּדֶשׁ וּבֵין קֹדֶשׁ הַקֳּדָשִׁים׃}
}
{
	\eP{
		And you shall place the pavilion under the clasps. And you
		shall bring the Ark of the Testimony there, inside the
		pavilion, and the pavilion will distinguish for you
		between the Holy and the Holy of Holies.
	}
}
{
%	\footnote{
%		under the clasps. The Greek text says “upon the frames,” meaning up
%		against them, or near them. In Hebrew, the frames are qršîm, and the
%		clasps are qrîm. The similarity of the two words led to the scribal
%		confusion here. The Greek fits with the understanding that the prket is
%		a pavilion, not a veil. The Masoretic Text makes little sense if it is a
%		pavilion. (If it were a veil, and if the Tabernacle were thirty cubits
%		long, then the Masoretic Text would mean that the veil hangs down under
%		the clasps, which come at the point of separation between the Holy place
%		and the Holy of Holies.) Footnote 26:33. bring the Ark of the Testimony
%		there, inside the pavilion. The ark is surrounded by a series of layers:
%		first the pavilion (the prket), then the Tabernacle (the linen tent),
%		then the tent (the goats’ hair tent), then the covering (the tent of
%		rams’ skins dyed red), then the covering of leather skins, then a
%		courtyard surrounded by linen hangings. The Israelites are thus given to
%		understand that there is a sequence of holiness, ever increasing, as one
%		comes closer to the sacred object in the center. Anyone who pursues
%		holiness must be aware that he or she must pass through a progression of
%		levels, with increasing awe and increasing danger as one comes closer to
%		the sacred. Footnote 26:33. the Holy. The space from the entrance up to
%		the prket is referred to as the Holy. Footnote 26:33. the Holy of
%		Holies. The space inside the prket is the Holy of Holies, meaning the
%		holiest of all holy places.
%	}
}

\Verse{34}
{
	\hP{וְנָתַתָּ אֶת הַכַּפֹּרֶת עַל אֲרוֹן הָעֵדֻת בְּקֹדֶשׁ הַקֳּדָשִׁים׃}
}
{
	\eP{
		And you shall place the atonement dais on the Ark of the
		Testimony in the Holy of Holies.
	}
}
{
}

\Verse{35}
{
	\hP{וְשַׂמְתָּ אֶת הַשֻּׁלְחָן מִחוּץ לַפָּרֹכֶת וְאֶת הַמְּנֹרָה נֹכַח הַשֻּׁלְחָן עַל צֶלַע הַמִּשְׁכָּן תֵּימָנָה וְהַשֻּׁלְחָן תִּתֵּן עַל צֶלַע צָפוֹן׃}
}
{
	\eP{
		And you shall place the table outside the pavilion, and
		the menorah opposite the table on the south side of the
		Tabernacle, and you shall place the table on the north
		side.
	}
}
{
}

\Verse{36}
{
	\hP{וְעָשִׂיתָ מָסָךְ לְפֶתַח הָאֹהֶל תְּכֵלֶת וְאַרְגָּמָן וְתוֹלַעַת שָׁנִי וְשֵׁשׁ משְׁזָר מַעֲשֵׂה רֹקֵם׃}
}
{
	\eP{
		“And you shall make a cover for the entrance of the tent:
		blue and purple and scarlet and woven linen—embroiderer’s
		work.
	}
}
{
}

\Verse{37}
{
	\hP{וְעָשִׂיתָ לַמָּסָךְ חֲמִשָּׁה עַמּוּדֵי שִׁטִּים וְצִפִּיתָ אֹתָם זָהָב וָוֵיהֶם זָהָב וְיָצַקְתָּ לָהֶם חֲמִשָּׁה אַדְנֵי נְחֹשֶׁת׃}
}
{
	\eP{
		And you shall make five acacia columns for the cover and
		plate them with gold, their hooks of gold, and cast five
		bases of bronze for them.
	}
}
{
}
